\documentclass{ieeeaccess}
\usepackage{cite}
\usepackage{amsmath,amssymb,amsfonts}
\usepackage{algorithmic}
\usepackage{graphicx}
\usepackage{textcomp}
\def\BibTeX{{\rm B\kern-.05em{\sc i\kern-.025em b}\kern-.08em
    T\kern-.1667em\lower.7ex\hbox{E}\kern-.125emX}}
\begin{document}
\history{Date of publication xxxx 00, 0000, date of current version xxxx 00, 0000.}
\doi{10.1109/TQE.2020.DOI}

\title{Circuit Optimization and Noise Robustness of Quantum Kernel Models for IoT Intrusion Detection}
% TODO: Replace with actual author names, affiliations, and IEEE memberships
\author{\uppercase{First A. Author}\authorrefmark{1}, \IEEEmembership{Fellow, IEEE},
\uppercase{Second B. Author\authorrefmark{2}, and Third C. Author,
Jr}.\authorrefmark{3},
\IEEEmembership{Member, IEEE}}
\address[1]{National Institute of Standards and 
Technology, Boulder, CO 80305 USA (email: author@boulder.nist.gov)}
\address[2]{Department of Physics, Colorado State University, Fort Collins, 
CO 80523 USA (email: author@lamar.colostate.edu)}
\address[3]{Electrical Engineering Department, University of Colorado, Boulder, CO 
80309 USA}
% TODO: Replace with actual funding/support acknowledgment or remove if not applicable
\tfootnote{This paragraph of the first footnote will contain support 
information, including sponsor and financial support acknowledgment. For 
example, ``This work was supported in part by the U.S. Department of 
Commerce under Grant BS123456.''}

\markboth
{Author \headeretal: Circuit Optimization and Noise Robustness of Quantum Kernel Models for IoT Intrusion Detection}
{Author \headeretal: Circuit Optimization and Noise Robustness of Quantum Kernel Models for IoT Intrusion Detection}

% TODO: Replace with actual corresponding author name and email
\corresp{Corresponding author: First A. Author (email: author@ boulder.nist.gov).}

\begin{abstract}
We examine how quantum circuit optimization affects kernel-based quantum machine learning for Internet of Things intrusion detection under ideal and noisy simulation. We evaluate three quantum kernel models, namely a quantum support vector classifier, a quantum voting ensemble, and a quantum weighted ensemble, on two binary intrusion-detection datasets: IoT Original Distribution and UNSW 2018 IoT Botnet Final 10 Best. The principal experiment uses 6 qubits and 5000 samples per dataset, compares Qiskit transpiler optimization levels 0 and 3, studies both full and linear entanglement, and evaluates 10 noise levels up to 5\% with 30 paired runs per condition. Because full 10-qubit noisy density-matrix evaluation at 5000 samples proved computationally impractical, we also conduct a supplementary 10-qubit extension at 2500 samples using the quantum support vector classifier and the quantum voting ensemble, full entanglement, four noise levels, and 10 runs. The results show that optimization effects are strongly conditional rather than uniform. The quantum voting ensemble is the most noise-robust model family across both datasets. On UNSW, full-entanglement quantum support vector and weighted-ensemble kernels can collapse under strong noise, whereas linear entanglement and the voting ensemble remain stable. On IoT, the voting ensemble remains strong even at the highest noise level. Holm-adjusted paired tests show statistically significant level-3 gains only in selected regimes, most notably for noisy weighted-ensemble configurations on UNSW. Overall, the findings show that robust quantum kernel design depends jointly on circuit optimization, entanglement topology, and model family, and cannot be inferred from depth reduction alone.
\end{abstract}

\begin{keywords}
Internet of Things intrusion detection, noise robustness, quantum circuit optimization, quantum kernel methods, quantum machine learning
\end{keywords}

\titlepgskip=-15pt

\maketitle

\section{Introduction}
\label{sec:introduction}

\PARstart{Q}{uantum} kernel methods are among the most plausible near-term quantum machine learning approaches because they separate quantum feature construction from classical decision boundaries. That separation is attractive for cybersecurity problems such as network intrusion detection, where the data are structured, high-dimensional, and often imbalanced. It does not, however, remove the central hardware constraint of the Noisy Intermediate-Scale Quantum era: kernel quality still depends on the circuits used to encode data, and those circuits remain vulnerable to depth growth, routing overhead, and accumulated two-qubit error.

For that reason, circuit optimization is not merely a compilation detail in quantum machine learning. Transpiler choices can alter circuit depth, two-qubit gate count, and ultimately the geometry of noisy quantum kernels experienced by the downstream classifier. Yet most prior work evaluates circuit optimization through proxy metrics such as depth or gate count, while most quantum machine learning studies report downstream accuracy without systematically comparing optimization strategies. The question that matters for practice is therefore end-to-end: when a quantum kernel model is trained on a real task under realistic simulated noise, when does optimization materially improve performance, and when does the model architecture itself dominate the outcome?

We address that question through a five-phase experimental framework built around paired comparisons between Qiskit transpiler level 0 (L0) and level 3 (L3). The completed study has two evidence tiers. The primary study is a fully repeated 6-qubit, 5000-sample experiment on both datasets with three model families, two entanglement topologies, 10 noise levels, and 30 runs per condition. The supplementary study is a reduced 10-qubit extension designed to preserve higher-qubit evidence after the full 10-qubit noisy design proved computationally infeasible. It uses 2500 samples, two model families, full entanglement only, four noise levels, and 10 runs per condition.

The paper makes four contributions. First, it provides an end-to-end evaluation of quantum circuit optimization for kernel-based intrusion detection, using downstream classification performance and runtime rather than circuit metrics alone. Second, it shows that optimization benefits depend on the model family, entanglement topology, and dataset; L3 helps materially in some noisy regimes, but it is not uniformly superior to L0. Third, it identifies the quantum voting ensemble as the most noise-robust model family in the study and shows that full entanglement can become brittle for some kernels under strong noise, especially on the UNSW dataset. Finally, it documents a practical limit of this workload: 10-qubit noisy kernel experiments are feasible, but only under a reduced supplementary design and at substantially greater computational cost than the primary 6-qubit study.

\section{Related Work}
\label{sec:related_work}

This section surveys the relevant literature across four intersecting research areas: quantum circuit optimization, quantum kernel methods and QML architectures, quantum ensemble learning, and QML applications to network intrusion detection. We conclude by identifying the specific gap that our work addresses.

\subsection{Quantum Circuit Optimization}
\label{sec:rw_optimization}

The NISQ era, characterized by Preskill~\cite{Preskill2018} as an epoch of devices with tens to hundreds of qubits operating without full error correction, imposes strict constraints on viable circuit depth. For a circuit with $g$ gates and per-gate error rate $\epsilon$, the probability of error-free execution degrades approximately as $(1-\epsilon)^g$~\cite{Cross2019}, making circuit optimization essential for any practical quantum application.

\subsubsection{Rule-Based Transpilation}
The most widely deployed optimization approach is rule-based transpilation, as implemented in Qiskit~\cite{Qiskit2024}, Cirq~\cite{Cirq2021}, and t$|$ket$\rangle$~\cite{Sivarajah2020}. Qiskit provides four optimization levels (0--3) with increasing aggressiveness: Level~0 performs only gate decomposition to the target basis gate set; Level~1 adds light optimization including inverse gate cancellation; Level~2 applies peephole optimization and commutative gate cancellation; and Level~3 applies the most aggressive pipeline including two-qubit block consolidation with KAK re-synthesis, unitary synthesis, and iterative optimization passes~\cite{Cross2022}. Dilillo et al.~\cite{Dilillo2023} systematically evaluated all four Qiskit optimization levels on 4,640 transpiled circuits and found that higher optimization levels generally improve circuit reliability, though the gains vary across circuit families. Nation et al.~\cite{Nation2025} presented a comprehensive benchmark of quantum computing software, comparing transpilation quality across frameworks under standardized conditions. Salm et al.~\cite{Salm2021} demonstrated that different compilers achieve superior results on different circuit classes, reinforcing the value of multi-method comparisons.

\subsubsection{ZX-Calculus Optimization}
ZX-calculus, introduced by Coecke and Duncan~\cite{Coecke2011}, provides a graphical language for quantum computation that represents circuits as ZX-diagrams and simplifies them via rewriting rules such as spider fusion, Hadamard complementarity, and local complementation~\cite{vandeWetering2020}. Kissinger and van de Wetering~\cite{Kissinger2020} demonstrated that ZX-calculus achieves optimal T-count reduction for Clifford+T circuits, and the PyZX library~\cite{Kissinger2019} has been shown to reduce circuit size by 20--40\% in many cases. Fischbach et al.~\cite{Fischbach2025} provided a comprehensive review of ZX-based circuit optimization, cataloging optimization strategies and identifying common limitations. However, ZX-calculus faces fundamental challenges with parameterized circuits: the extraction of gate sequences from simplified diagrams is computationally hard~\cite{deBeaudrap2020}, and optimization may not preserve the symbolic parameter structure required for variational and kernel-based QML circuits. Van de Wetering et al.~\cite{vandeWetering2025} recently addressed this limitation, proving that the PyZX optimization algorithm is optimal for compiling parameterized circuits under certain conditions, though practical extraction of parameterized gates from ZX-diagrams remains an open challenge.

\subsubsection{Alternative Compiler Frameworks}
Cambridge Quantum Computing's t$|$ket$\rangle$~\cite{Sivarajah2020} provides distinct optimization strategies emphasizing phase gadget compilation and peephole optimization. Nam et al.~\cite{Nam2018} demonstrated that aggressive optimization can reduce two-qubit gate counts by up to 50\% for certain circuit families. Liu et al.~\cite{Liu2021} proposed relaxed peephole optimization, achieving better gate reduction than Qiskit's most aggressive level for specific circuit patterns. Yan et al.~\cite{Yan2024} provided a comprehensive survey of quantum circuit synthesis and compilation with over 60 citations, unifying rule-based, algebraic, and synthesis-based approaches under a common framework and identifying complementary strengths among different methods.

\subsubsection{Machine Learning Approaches}
Recent work has applied machine learning to circuit optimization. F{\"o}sel et al.~\cite{Fosel2021} pioneered deep reinforcement learning for circuit optimization, training a convolutional neural network agent to reduce circuit depth by 27\% and gate count by 15\% on 12-qubit random circuits. Paler et al.~\cite{Paler2023} addressed qubit mapping optimization using machine learning, introducing QXX-MLP to reduce SWAP overhead on connectivity-constrained hardware. Ruiz et al.~\cite{Ruiz2025} developed AlphaTensor-Quantum, a deep reinforcement learning method published in Nature Machine Intelligence that exploits the relationship between T-count optimization and tensor decomposition, discovering efficient algorithms for circuits in Shor's algorithm and quantum chemistry. Kremer et al.~\cite{Kremer2024} demonstrated practical reinforcement learning for quantum circuit synthesis and transpiling with improved compilation quality. Wang et al.~\cite{Wang2022QuantumNAS} proposed QuantumNAS, a noise-adaptive search framework for finding robust quantum circuits that explicitly accounts for hardware noise during architecture search.

\subsection{Quantum Kernel Methods}
\label{sec:rw_kernels}

Quantum kernel methods represent one of the most promising near-term QML approaches. Schuld and Killoran~\cite{Schuld2019} formulated supervised learning in quantum feature Hilbert spaces, while Havl\'{i}\v{c}ek et al.~\cite{Havlicek2019} demonstrated supervised learning with quantum-enhanced feature spaces on IBM hardware. These methods encode classical data into quantum states via feature maps and compute kernel functions as the squared overlap $k(\mathbf{x}_i, \mathbf{x}_j) = |\langle\phi(\mathbf{x}_i)|\phi(\mathbf{x}_j)\rangle|^2$, which can then be used with classical SVMs.

Huang et al.~\cite{Huang2021} provided theoretical foundations for quantum kernel advantage, demonstrating that quantum kernels can achieve exponential separations from classical kernels for certain data distributions while also identifying the ``exponential concentration'' phenomenon, in which feature maps may produce uninformative kernels for high-dimensional data. Sim et al.~\cite{Sim2019} systematically characterized the expressibility and entangling capability of parameterized circuits, establishing quantitative metrics for comparing feature map designs. The choice of feature map, including popular designs such as ZZ~\cite{Havlicek2019} and Pauli feature maps~\cite{Schuld2019}, directly impacts both expressibility and circuit depth, making feature map selection a primary target for optimization studies.

From a variational perspective, McClean et al.~\cite{McClean2018} identified barren plateaus in quantum neural network training landscapes, where gradients vanish exponentially with system size. Holmes et al.~\cite{Holmes2022} further connected ansatz expressibility to gradient magnitudes, demonstrating that highly expressive circuits are more prone to barren plateaus. These findings underscore the tension between circuit complexity and trainability, motivating optimization strategies that maintain model expressivity while minimizing noise-susceptible circuit depth.

\subsection{Quantum Ensemble Methods}
\label{sec:rw_ensembles}

Quantum ensemble methods combine multiple quantum models to improve prediction robustness. Schuld and Petruccione~\cite{Schuld2018Ensemble} demonstrated that quantum ensembles of classifiers can achieve improved generalization over individual quantum classifiers. Qin et al.~\cite{Qin2022} proposed improving quantum classifier performance on NISQ computers through classical voting strategies from ensemble learning, training quantum classifiers across multiple quantum backends to mitigate hardware-specific noise effects.

Tolotti et al.~\cite{Tolotti2023} investigated ensembles of quantum classifiers, analyzing how combining heterogeneous quantum models via majority voting improves classification stability under noise. Incudini et al.~\cite{Incudini2023} demonstrated resource savings through ensemble techniques for quantum neural networks, showing that combining smaller circuits can match the performance of larger individual circuits at lower noise cost. Li et al.~\cite{Li2024ensemble} proposed ensemble-learning error mitigation for variational quantum shallow-circuit classifiers, using ensemble diversity to address noise without explicit error mitigation protocols. These results collectively suggest that ensemble strategies provide a practical pathway to noise-resilient quantum classification on NISQ hardware.

However, prior quantum ensemble approaches derive diversity from classical mechanisms such as data sampling, model architecture variation, or multi-backend execution, without exploiting the quantum-specific capability of encoding identical data into geometrically distinct Hilbert space regions. In our prior work~\cite{AuthorQVEQWE2025}, we addressed this gap by introducing two novel quantum ensemble methods, Quantum Voting Ensemble (QVE) and Quantum Weighted Ensemble (QWE), which achieve diversity at the quantum representation level by combining classifiers with heterogeneous feature encodings. QVE pairs Z and ZZ feature maps through hard majority voting, while QWE combines feature maps with adaptive weight assignment. Evaluation on the IoTID20 dataset demonstrated that QVE achieves 99.5\% accuracy, the highest among all quantum configurations evaluated, and that the 4.2 to 7.8 percentage point accuracy gap over Quantum Random Forest validates encoding-level diversity as a superior alternative to bootstrap aggregation with homogeneous circuits.

\subsection{QML for Network Intrusion Detection}
\label{sec:rw_nids}

Network intrusion detection represents a well-studied domain for classical machine learning~\cite{Buczak2016,Aldweesh2020}, with modern approaches employing deep learning~\cite{Vinayakumar2019} and ensemble methods. IoT networks present particular challenges due to high-dimensional feature spaces and evolving attack patterns~\cite{Cook2019}.

The application of QML to cybersecurity has gained momentum in recent years. Kalinin and Krundyshev~\cite{Kalinin2023} explored security intrusion detection using quantum machine learning techniques, evaluating both QSVM and quantum convolutional neural networks and demonstrating comparable accuracy to classical methods while identifying network intrusion detection as a candidate for quantum speedup. Kumar and Swarnkar~\cite{Kumar2025} proposed QuIDS, a quantum support vector machine-based intrusion detection system specifically designed for IoT networks, using 4-qubit quantum feature maps to classify network traffic. Alomari and Kumar~\cite{Alomari2023} introduced DEQSVC, combining dimensionality reduction and encoding techniques within a quantum support vector classifier for DDoS attack detection, executed on IBM quantum hardware. Elsedimy et al.~\cite{Elsedimy2024} proposed a hybrid quantum support vector machine with improved Grey Wolf optimizer for intrusion detection, achieving competitive detection rates. Cultice et al.~\cite{Cultice2025} recently investigated quantum-hybrid SVMs for anomaly detection in industrial control systems, comparing quantum kernel performance against traditional kernel methods across multiple CPS datasets.

The IoTID20 dataset~\cite{ullah2020scheme} used in this work has become a standard benchmark for IoT intrusion detection research~\cite{Houkan2024,Adefemi2025}. However, prior QML studies on network intrusion detection have focused primarily on classification accuracy without systematically studying how circuit optimization affects model performance under noise, which is the central question addressed here.

\subsection{Gap in the Literature}
\label{sec:rw_gap}

Despite extensive work on circuit optimization and QML individually, systematic studies of how optimization affects end-to-end QML performance remain scarce. Most circuit-optimization studies evaluate abstract benchmarks such as gate count and circuit depth rather than downstream task performance~\cite{Yan2024,Fischbach2025}. By contrast, QML studies typically employ default compiler settings without comparative analysis of optimization strategies~\cite{Kalinin2023,Kumar2025}. The interaction between optimization level and noise resilience has received limited attention. Existing work suggests that resilience gains can be substantial~\cite{Murali2020}, but no prior study has evaluated this question systematically across multiple models, feature maps, noise levels, and entanglement topologies with rigorous statistical testing.

Furthermore, no prior study has examined the practical limitations of applying alternative optimization paradigms such as ZX-calculus and Pytket to parameterized QML circuits, where preserving symbolic parameter structure is essential. While our prior work~\cite{AuthorQVEQWE2025} established the effectiveness of QVE and QWE for IoT intrusion detection under noiseless simulation, it did not investigate how circuit optimization alters ensemble performance under realistic noise conditions, which is the central question addressed here.

To address these gaps, we adopt a five-phase experimental framework that compares multiple optimization methodologies on QML circuits, evaluates end-to-end classification performance rather than proxy metrics alone, performs a fully repeated primary noise study at 6 qubits with 30 paired runs per configuration together with a reduced supplementary 10-qubit extension, and examines three distinct QML architectures, QSVC, QVE, and QWE, to determine whether optimization and entanglement effects extend beyond a single kernel model.


%% ====================================================================
%% METHODOLOGY SECTION — Circuit Optimization Experiment
%% ====================================================================

\section{Experimental Methodology}
\label{sec:methodology}

This section describes the experimental framework used to evaluate how quantum circuit optimization changes the performance of quantum kernel models for IoT intrusion detection under ideal and noisy simulation. The central question is \emph{when circuit optimization materially improves accuracy, robustness, and runtime for quantum kernel classifiers}. The final design consists of a primary 6-qubit study and a reduced supplementary 10-qubit extension. The primary study uses 30 paired runs per condition, whereas the 10-qubit extension uses 10 runs per condition and is interpreted as supplementary evidence rather than as a direct replacement for the primary factorial comparison.

\subsection{Dataset and Preprocessing}
\label{sec:dataset}

We evaluate two binary intrusion-detection datasets. The first is the IoT Network Intrusion Dataset (IoTID20)~\cite{ullah2020scheme}, from which we use a cleaned tabular version denoted \emph{IoT Original Distribution}. The second is \emph{UNSW 2018 IoT Botnet Final 10 Best}, a compact intrusion-detection dataset with a preselected feature subset. Using both datasets allows us to test whether optimization and noise effects replicate across two different IoT intrusion-detection distributions rather than overfitting conclusions to a single benchmark.

\subsubsection{Data Wrangling}
Prior to the quantum experiments, we applied a systematic data wrangling pipeline to prepare a clean experimental dataset from the raw IoTID20 data:
\begin{enumerate}
\item \textbf{Removal of non-predictive features:} Four columns that are either unique identifiers or sources of data leakage were removed: \texttt{Flow\_ID} (unique flow identifier), \texttt{Src\_IP} and \texttt{Dst\_IP} (IP addresses), and \texttt{Timestamp} (temporal information).
\item \textbf{Missing value imputation:} Numeric columns were imputed with the column median, and categorical columns with the mode.
\item \textbf{Infinite value handling:} Infinite values ($\pm\infty$) were replaced with NaN and subsequently imputed with the column median.
\item \textbf{Duplicate removal:} Exact duplicate rows were removed, reducing the dataset from 625,783 to 293,696 instances, a 53\% reduction that indicates substantial redundancy in the raw network traffic captures.
\item \textbf{Outlier capping:} Outliers were capped using the interquartile range (IQR) method: values below $Q_1 - 1.5 \cdot \text{IQR}$ or above $Q_3 + 1.5 \cdot \text{IQR}$ were clipped to the respective bounds, preserving all instances while mitigating extreme values.
\item \textbf{Feature engineering:} Three domain-informed ratio features were created: forward-to-backward packet ratio (\texttt{Fwd\_Bwd\_Pkt\_Ratio}), forward-to-backward payload length ratio (\texttt{Fwd\_Bwd\_Len\_Ratio}), and overall packet rate (\texttt{Pkt\_Rate}), computed from flow duration and packet counts.
\item \textbf{Categorical encoding:} Remaining categorical features were transformed via label encoding.
\end{enumerate}

The resulting IoT dataset, denoted \emph{IoT Original Distribution}, contains 293,696 instances across 85 columns: 82 feature columns and three target columns (\texttt{Label}, \texttt{Cat}, \texttt{Sub\_Cat}). The binary class distribution is preserved from the raw data, with 266,166 anomaly instances (90.6\%) and 27,530 normal instances (9.4\%), yielding an imbalance ratio of approximately 9.7:1.

For the UNSW dataset, we use the curated file provided in the experimental repository. Because that file already contains a compact selected-feature representation, the preprocessing pipeline focuses on binary label extraction, removal of identifier or leakage-prone columns when present, stratified sampling, feature scaling to $[0, \pi]$, and projection to the target qubit dimension.

\subsubsection{Binary Formulation}
We cast the problem as binary classification (Normal vs.\ Anomaly) by grouping all attack categories into a single anomaly class. This formulation aligns with practical deployment scenarios where the primary objective is distinguishing benign from malicious traffic regardless of attack type.

\subsubsection{Experiment-Specific Preprocessing}
During each experimental run, the selected dataset undergoes additional preprocessing tailored to the quantum circuit requirements. Categorical columns such as \texttt{Cat} and \texttt{Sub\_Cat} are removed whenever they directly encode attack-type labels, thereby preventing label leakage. Binary labels are then assigned with \texttt{LabelEncoder}, mapping Normal to 0 and Anomaly to 1. The dataset is downsampled to $N$ samples through stratified sampling so that the original class distribution is preserved. Any residual NaN values are replaced with 0.0, positive infinite values are clipped to $\pi$, and negative infinite values are set to 0.0. When the number of available numerical features exceeds $2n$, where $n$ is the qubit count, \texttt{SelectKBest} with mutual information (\texttt{mutual\_info\_classif}) retains the $2n$ most informative features. The retained features are then mapped to the interval $[0, \pi]$ with \texttt{MinMaxScaler}, matching the domain of the quantum rotation gates while preserving angular separation. Finally, Principal Component Analysis (PCA) reduces the feature space to exactly $n$ components, producing a one-to-one correspondence between features and qubits.

Data is partitioned using a 70:30 stratified train-test split. The primary 6-qubit study uses 30 distinct random seeds, while the supplementary 10-qubit extension uses 10 distinct seeds.

\subsubsection{Experimental Configurations}
The completed experiment has two configuration tiers. The primary study uses both datasets at 6 qubits with 5000 samples and serves as the main controlled comparison for inferential claims. The supplementary extension uses both datasets at 10 qubits with 2500 samples. This reduced design was adopted after the full 10-qubit, 5000-sample noisy kernel workload proved computationally impractical under density-matrix simulation. The extension preserves higher-qubit evidence while narrowing the scope of the claim.

\subsection{Quantum Circuit Architecture}
\label{sec:circuit_architecture}

\subsubsection{Feature Maps}
We employ three quantum feature maps that correspond directly to the trained models in the experiment.

\paragraph{Z Feature Map} The Z feature map applies single-qubit rotations encoding each feature $x_i$ directly into quantum phase:
\begin{equation}
U_Z(\mathbf{x}) = \bigotimes_{i=1}^{n} R_Z(2x_i)\, H\, |0\rangle
\label{eq:z_map}
\end{equation}
where $R_Z(\theta) = e^{-i\theta Z/2}$ is the Pauli-$Z$ rotation gate and $H$ is the Hadamard gate. This encoding produces separable product states with circuit depth $\mathcal{O}(1)$, offering inherent noise resilience due to the absence of entangling gates.

\paragraph{ZZ Feature Map} The ZZ feature map augments single-qubit rotations with two-qubit entangling gates that encode pairwise feature interactions:
\begin{equation}
U_{ZZ}(\mathbf{x}) = \exp\!\left(i \sum_{i<j} \phi_{ij}(\mathbf{x})\, Z_i Z_j\right) U_Z(\mathbf{x})
\label{eq:zz_map}
\end{equation}
where $\phi_{ij}(\mathbf{x}) = (\pi - x_i)(\pi - x_j)$ captures second-order feature correlations through entangled quantum states. This map serves as the primary feature encoding for all single-model experiments with $r = 2$ repetition layers and full entanglement connectivity.

\paragraph{Pauli Feature Map} Configured with Pauli operators $\{Z, ZZ\}$, this encoding provides an alternative implementation of ZZ-style feature interactions. It is used as the secondary encoding in the Quantum Weighted Ensemble (QWE).

\subsubsection{Entanglement Topologies}
The final study evaluates two entanglement patterns that determine the connectivity structure of two-qubit gates. Linear entanglement uses nearest-neighbor $\text{CX}$ gates ($i \to i{+}1$), which closely matches the native topology of many superconducting processors. Full entanglement uses all-to-all $\text{CX}$ connectivity and maximizes interaction at the cost of greater circuit depth. The primary 6-qubit study includes both patterns, whereas the supplementary 10-qubit extension retains only full entanglement to control runtime.

\subsubsection{Quantum Kernel}
Each encoded quantum state defines a fidelity-based quantum kernel measuring similarity between data points in Hilbert space:
\begin{equation}
K(\mathbf{x}_i, \mathbf{x}_j) = \bigl|\langle \psi(\mathbf{x}_i) | \psi(\mathbf{x}_j) \rangle\bigr|^2
\label{eq:kernel}
\end{equation}

For ideal (noiseless) evaluation, we compute exact statevectors $|\psi(\mathbf{x})\rangle$ via the Qiskit Aer \texttt{statevector} simulator and evaluate the squared inner product using GPU-accelerated matrix multiplication. For noisy evaluation, we employ a sampling-based kernel using Hellinger fidelity (Section~\ref{sec:noise_model}).

\subsection{Circuit Optimization Framework}
\label{sec:optimization}

The central methodological comparison in this work is between two Qiskit transpiler optimization levels. Level 0 (L0) performs only gate decomposition to the target basis gate set $\{u, cx, rz, sx, x\}$ and omits circuit simplification, gate cancellation, and routing optimization. It therefore serves as the baseline representation of the circuit as constructed. Level 3 (L3) applies Qiskit's most aggressive optimization pipeline, including gate cancellation and merging, two-qubit block consolidation with KAK re-synthesis, commutative gate cancellation, unitary synthesis, layout and routing optimization, and iterative optimization passes.

Both optimization levels target a linear coupling map that models realistic nearest-neighbor qubit connectivity on current superconducting hardware. All transpilations use a fixed seed (\texttt{seed\_transpiler}$= 42$) for reproducibility.

\subsubsection{Circuit Metrics}
For each circuit configuration, we record circuit depth, total gate count, two-qubit gate count, and a NISQ feasibility score. Circuit depth is the length of the longest path through the circuit's directed acyclic graph. Total gate count measures the number of gates after decomposition to the basis gate set, while two-qubit gate count tracks the number of $\text{CX}$ gates, which dominate the error budget on NISQ hardware. The NISQ feasibility score is a composite measure on a 0--100 scale that incorporates estimated circuit fidelity, coherence-budget utilization, and gate-error accumulation using IBM Falcon r5.11 device parameters.

The estimated circuit fidelity under the independent error model is:
\begin{equation}
\mathcal{F}_{\text{circuit}} = (1 - \epsilon_{1q})^{g_{1q}} \cdot (1 - \epsilon_{2q})^{g_{2q}} \cdot (1 - \epsilon_{\text{ro}})
\label{eq:fidelity}
\end{equation}
where $\epsilon_{1q}$, $\epsilon_{2q}$, and $\epsilon_{\text{ro}}$ are single-qubit, two-qubit, and readout error rates, respectively, and $g_{1q}$, $g_{2q}$ are the corresponding gate counts.

\subsubsection{Multi-Method Optimization Analysis}
In addition to the Qiskit transpiler, we evaluate two alternative optimization methodologies in a supporting phase to assess whether fundamentally different paradigms yield superior circuit reductions for QML feature maps. The first is ZX-calculus via \texttt{pyzx}, a graph-theoretic simplification framework that represents quantum circuits as ZX-diagrams and applies rewriting rules such as spider fusion, Hadamard complementarity, and local complementation to reduce circuit complexity; in this study, we use the \texttt{teleport\_reduce} strategy to improve circuit extraction. The second is the Pytket compiler, which applies a combination of peephole simplification, redundancy removal, and multi-qubit commutation passes. In practice, neither ZX-calculus nor Pytket could serve as a drop-in replacement for the parameterized model-training circuits because the symbolic structure required by the QML feature maps was not preserved by the workable optimization pipeline. Consequently, all end-to-end model comparisons in the final study rely on Qiskit L0 and L3 circuits, while the alternative optimizers are treated as supporting evidence on the limits of optimizing parameterized QML circuits.

\subsection{QML Models}
\label{sec:models}

We evaluate the impact of circuit optimization on three QML model architectures spanning kernel-based and ensemble paradigms. The ensemble models (QVE and QWE) were introduced in our prior work~\cite{AuthorQVEQWE2025}; here we summarize their design and investigate how circuit optimization affects their performance under noise.

\subsubsection{Quantum Support Vector Classifier (QSVC)}
The primary model computes a precomputed kernel matrix $K \in \mathbb{R}^{N \times N}$ using the quantum fidelity kernel~\eqref{eq:kernel}. Classification is then performed with a GPU-accelerated cuML \texttt{SVC} configured with a precomputed kernel, $C = 1.0$, balanced class weights, a maximum of 10{,}000 iterations, and an 8192 MB kernel cache. This precomputed-kernel formulation decouples quantum circuit evaluation from classical optimization and enables efficient GPU utilization.

\subsubsection{Quantum Voting Ensemble (QVE)}
QVE achieves ensemble diversity through heterogeneous quantum feature encodings rather than data sampling. Two QSVC base classifiers are constructed using identical training data but geometrically distinct feature maps: $h_Z$ employing the Z feature map and $h_{ZZ}$ employing the ZZ feature map. The ensemble prediction is obtained via hard majority voting:
\begin{equation}
\hat{y}_{\text{QVE}} = \arg\max_{c \in \mathcal{Y}} \sum_{m=1}^{M} \mathbf{1}[h_m(\mathbf{x}) = c]
\label{eq:qve}
\end{equation}
where $M = 2$ and $\mathbf{1}[\cdot]$ is the indicator function.

\subsubsection{Quantum Weighted Ensemble (QWE)}
QWE extends hard voting with soft voting via probability averaging. Each base classifier produces class probability estimates, and the ensemble averages these probabilities before selecting the class with the highest mean probability:
\begin{equation}
\hat{y}_{\text{QWE}} = \arg\max_{c \in \mathcal{Y}} \frac{1}{M} \sum_{m=1}^{M} P_m(c \mid \mathbf{x})
\label{eq:qwe}
\end{equation}
where $P_m(c \mid \mathbf{x})$ is the class probability estimate from the $m$-th classifier. QWE pairs the ZZ feature map with the Pauli feature map, providing a different encoding combination than QVE.

\subsection{Noise Model and Noisy Kernel Evaluation}
\label{sec:noise_model}

To assess circuit optimization benefits under realistic NISQ conditions, we construct a parameterized noise model combining depolarizing and thermal relaxation channels.

\subsubsection{Noise Channel Construction}
For single-qubit gates, we compose a depolarizing channel with probability $p_{1q}$ and a thermal relaxation channel parameterized by coherence times $T_1$ and $T_2$ with gate duration $\tau_{1q}$:
\begin{equation}
\mathcal{E}_{1q} = \mathcal{E}_{\text{depol}}(p_{1q}) \circ \mathcal{E}_{\text{thermal}}(T_1, T_2, \tau_{1q})
\label{eq:noise_1q}
\end{equation}
For two-qubit gates, a depolarizing channel with probability $p_{2q}$ is applied:
\begin{equation}
\mathcal{E}_{2q} = \mathcal{E}_{\text{depol}}(p_{2q})
\label{eq:noise_2q}
\end{equation}

The base noise parameters model a conservative superconducting device: $p_{1q} = 10^{-3}$, $p_{2q} = 10^{-2}$, $T_1 = 50~\mu$s, $T_2 = 70~\mu$s, $\tau_{1q} = 50$ ns, $\tau_{2q} = 300$ ns. The two-qubit error rate is maintained at $10\times$ the single-qubit rate across all noise scaling levels, reflecting realistic hardware error ratios.

\subsubsection{Noise-Scaled Experiment}
The primary 6-qubit study sweeps 10 noise levels by scaling the base single-qubit error rate:
\begin{align}
p_{1q} \in \{&0,\, 3{\times}10^{-4},\, 5{\times}10^{-4},\, 10^{-3},\, 2{\times}10^{-3}, \notag \\
&5{\times}10^{-3},\, 8{\times}10^{-3},\, 10^{-2},\, 2{\times}10^{-2},\, 5{\times}10^{-2}\}. \notag
\end{align}
The supplementary 10-qubit extension uses a reduced four-level grid,
\begin{equation}
p_{1q} \in \{0,\, 2{\times}10^{-3},\, 10^{-2},\, 5{\times}10^{-2}\},
\end{equation}
chosen to retain ideal, low-noise, moderate-noise, and stress-test conditions while keeping the higher-qubit density-matrix workload tractable. The primary study uses 30 independent runs per condition, and the 10-qubit extension uses 10 runs per condition.

\subsubsection{Noisy Kernel via Hellinger Fidelity}
Under noise, the statevector-based kernel~\eqref{eq:kernel} is replaced by a sampling-based kernel. Each data point $\mathbf{x}$ is encoded, the noisy circuit is simulated using the density matrix method (\texttt{AerSimulator} with \texttt{method=`density\_matrix'}), and measurement outcomes are sampled with 1024 shots to obtain a probability distribution $\mathbf{p}(\mathbf{x})$. The noisy kernel is then computed as the Hellinger fidelity (Bhattacharyya coefficient):
\begin{equation}
K_{\text{noisy}}(\mathbf{x}_i, \mathbf{x}_j) = \sum_{k} \sqrt{p_k(\mathbf{x}_i) \cdot p_k(\mathbf{x}_j)}
\label{eq:hellinger}
\end{equation}
which approximates the ideal fidelity kernel in the zero-noise, infinite-shot limit. The full noisy kernel matrix is computed as $K = \sqrt{P_1}\, (\sqrt{P_2})^T$ where rows of $P$ are the measurement probability vectors.

\subsection{Experiment Design}
\label{sec:experiment_design}

The experiment comprises five phases, each targeting a specific research question.

\subsubsection{Phase 1: Multi-Method Optimization Comparison}
Compares circuit metrics across optimization methods and evaluates a supporting QSVC L0 versus L3 sanity check under ideal simulation.

\subsubsection{Phase 2: Unoptimized vs.\ Optimized Comparison}
This phase extends the ideal-simulation comparison across feature-map variants in order to determine whether optimization effects depend on the choice of feature map before noise is introduced.

\subsubsection{Phase 3: Qiskit Optimization Levels (0--3)}
Evaluates circuit metrics and QSVC accuracy across Qiskit optimization levels 0--3 to characterize the marginal benefit of increasingly aggressive transpilation.

\subsubsection{Phase 4: Comprehensive Model-Noise-Optimization Analysis}
Phase 4 is the core paper experiment. In the primary study, it evaluates all combinations of 3 models (QSVC, QVE, QWE), 2 optimization levels (L0, L3), 2 entanglement topologies (full, linear), 10 noise levels, and 30 runs, yielding 3600 individual model evaluations per dataset. In the supplementary 10-qubit extension, Phase 4 is restricted to 2 models (QSVC, QVE), 2 optimization levels, full entanglement only, 4 noise levels, and 10 runs, yielding 160 model evaluations per dataset. For each configuration, the ideal kernel uses exact statevector fidelity and the noisy kernels use Hellinger fidelity derived from noisy density-matrix simulation.

\subsubsection{Phase 5: Cross-Validation with Statistical Testing}
Performs leakage-safe 5-fold stratified cross-validation for supporting generalization estimates. Cross-validation is reported separately from the noisy Phase 4 paired comparisons and is used to verify that the ideal kernel models remain stable across repeated folds and seeds.

\subsection{Statistical Analysis}
\label{sec:statistics}

The primary study uses 30 paired runs per condition, satisfying conventional sample-size guidance for parametric paired inference. The supplementary 10-qubit extension uses 10 paired runs per condition and is therefore interpreted more cautiously. Statistical analysis combines paired $t$-tests (\texttt{scipy.stats.ttest\_rel}) for L0 versus L3 accuracy distributions, Holm correction to control the family-wise error rate within each result family, paired Cohen's $d$ to quantify effect magnitude, and 95\% confidence intervals computed as $\bar{x} \pm t_{\alpha/2} \cdot s / \sqrt{n}$. We also report descriptive noise-response patterns, including accuracy degradation and runtime inflation from the ideal to the highest-noise conditions.

Evaluation metrics per model include accuracy, balanced accuracy, weighted and macro F1 scores, precision, recall, specificity, Matthews Correlation Coefficient (MCC), Cohen's kappa, and ROC-AUC.

\subsection{Computational Infrastructure}
\label{sec:infrastructure}

All experiments are executed on a system equipped with three NVIDIA RTX A6000 GPUs (48~GB VRAM each). Ideal kernel evaluation uses Qiskit Aer statevector simulation with GPU acceleration. Noisy kernel evaluation uses density-matrix simulation with GPU-backed linear algebra and blockwise kernel construction to keep higher-qubit workloads within memory limits.

Kernel matrix computation is parallelized across the available GPUs, followed by GPU-accelerated fidelity computation via CuPy. Classical SVM training uses RAPIDS cuML with precomputed kernel matrices resident in GPU memory. The substantial runtime gap between the 6-qubit and 10-qubit noisy workloads is therefore a property of the simulation regime itself rather than of an unoptimized software path.

\section{Results}
\label{sec:results}

The results synthesized in this section are drawn from the complete archive produced by the two completed experiment batches: the primary 6-qubit, 5000-sample batch and the supplementary 10-qubit, 2500-sample batch. Together, these batches produced 52 result artifacts, with 26 files per batch and 13 files per dataset-specific configuration. Rather than reproducing every machine-readable file verbatim, the manuscript integrates the full archive by result family: circuit-level diagnostics, optimizer-comparison outputs, comprehensive noise-study outputs, cross-validation summaries, and batch-level summaries. Table~\ref{tab:result_archive} makes this coverage explicit so that all five experimental phases are represented in the paper narrative.

\begin{table*}[t]
\caption{Coverage of the complete result archive across the two completed experiment batches. The 6-qubit and 10-qubit batches each contain dataset-specific result bundles for IoT and UNSW, giving 52 archived artifacts in total.}
\label{tab:result_archive}
\centering
\setlength{\tabcolsep}{4.5pt}
\resizebox{0.98\textwidth}{!}{%
\begin{tabular}{|l|l|l|c|}
\hline
Result family & Files included & Experimental role & Total files \\
\hline
Circuit-level diagnostics & \texttt{circuit\_metrics}, \texttt{multi\_method\_optimization}, \texttt{optimization\_levels} & Phases 1--3 circuit compression, gate-count, and transpilation-level analysis & 12 \\
Model-level optimizer comparisons & \texttt{optimization\_comparison}, \texttt{multi\_method\_model\_results}, \texttt{optimization\_statistics}, \texttt{multi\_method\_statistics} & Phase 1--3 ideal-simulation accuracy comparisons and paired statistical summaries & 16 \\
Comprehensive noisy-study outputs & \texttt{comprehensive\_model\_noise}, \texttt{comprehensive\_model\_noise\_runs}, \texttt{comprehensive\_model\_noise\_stats} & Phase 4 per-configuration means, per-run distributions, and L0/L3 statistical contrasts under noise & 12 \\
Cross-validation outputs & \texttt{cv\_results} & Phase 5 leakage-safe 5-fold generalization estimates across qubit scales and datasets & 4 \\
Batch-level summaries & \texttt{experiment\_summary}, \texttt{journal\_summary} & Configuration metadata, runtime rollups, and publication-oriented batch summaries & 8 \\
\hline
Total archived results & All families above & Full coverage of the completed 6-qubit and 10-qubit evidence tiers & 52 \\
\hline
\end{tabular}%
}
\end{table*}

\begin{table*}[t]
\caption{Primary 6-qubit Phase 4 summary at the ideal setting and at the highest noise level ($p_{1q}=0.05$). IoT and UNSW denote the two intrusion-detection datasets used in the study.}
\label{tab:primary6q_summary}
\centering
\setlength{\tabcolsep}{5pt}
\resizebox{0.98\textwidth}{!}{%
\begin{tabular}{|l|l|c|l|c|}
\hline
Dataset & Best ideal configuration & Accuracy (\%) & Best highest-noise configuration & Accuracy (\%) \\
\hline
IoT & QWE, full, L3 & 95.62 & QVE, full, L3 & 96.78 \\
UNSW & QWE, full, L0/L3 & 99.62 & QWE, linear, L0/L3 & 98.91 \\
\hline
\end{tabular}%
}
\end{table*}

\begin{table*}[t]
\caption{Selected Holm-adjusted paired L0 versus L3 accuracy contrasts from the noisy-study statistics files. Positive $\Delta$ favors L3, whereas negative $\Delta$ favors L0.}
\label{tab:holm_selected}
\centering
\setlength{\tabcolsep}{4.5pt}
\resizebox{0.98\textwidth}{!}{%
\begin{tabular}{|l|l|l|c|c|c|}
\hline
Study block & Model / entanglement & Noise level & L0 acc. (\%) & L3 acc. (\%) & $\Delta$L3$-$L0 (pp) \\
\hline
IoT 6q primary & QWE / full & 0.02 & 91.24 & 93.81 & +2.57 \\
IoT 6q primary & QWE / full & 0.01 & 94.06 & 94.47 & +0.41 \\
IoT 6q primary & QVE / full & 0.02 & 96.32 & 95.47 & -0.85 \\
UNSW 6q primary & QWE / full & 0.05 & 9.53 & 74.88 & +65.35 \\
UNSW 6q primary & QWE / full & 0.02 & 98.08 & 98.91 & +0.83 \\
UNSW 10q supplementary & QSVC / full & 0.002 & 82.04 & 93.29 & +11.25 \\
UNSW 10q supplementary & QVE / full & 0.05 & 96.52 & 97.01 & +0.49 \\
UNSW 10q supplementary & QVE / full & 0.01 & 97.68 & 97.08 & -0.60 \\
\hline
\end{tabular}%
}
\end{table*}

\begin{table*}[t]
\caption{Supplementary 10-qubit robustness and runtime summary. Reported values give the best observed ideal and highest-noise accuracies within each model family across L0/L3, together with the largest runtime inflation from the ideal to the highest-noise condition. Dataset abbreviations follow Table~\ref{tab:primary6q_summary}.}
\label{tab:tenq_summary}
\centering
\setlength{\tabcolsep}{5pt}
\resizebox{0.98\textwidth}{!}{%
\begin{tabular}{|l|l|c|c|c|c|}
\hline
Dataset & Model family & Best ideal acc. (\%) & Best high-noise acc. (\%) & Accuracy drop (pp) & Max runtime inflation \\
\hline
IoT & QSVC & 95.99 & 90.67 & 5.32 & $5.79\times$ \\
IoT & QVE & 97.57 & 97.25 & 0.32 & $6.84\times$ \\
UNSW & QSVC & 97.01 & 90.00 & 7.01 & $5.68\times$ \\
UNSW & QVE & 98.05 & 97.01 & 1.04 & $6.91\times$ \\
\hline
\end{tabular}%
}
\end{table*}

\subsection{Supporting Optimization Findings}

The supporting phases establish two boundary conditions for the main study. First, aggressive Qiskit transpilation reduces circuit metrics for the feature maps considered here, but the downstream performance effect cannot be inferred from gate counts alone. Table~\ref{tab:circuit_metrics} quantifies the circuit-level impact of L3 optimization across both qubit scales. At 6 qubits, L3 reduces two-qubit (CX) gate count by 5.3\% for ZZ $r{=}2$ and by 11.1\% for ZZ $r{=}3$, while depth is unchanged or marginally increased. At 10 qubits, the reductions are substantially larger: depth drops by 17.7\%, CX gates by 21.5\% for ZZ $r{=}2$, and by 16.7\% for ZZ $r{=}3$. These asymmetric reductions across qubit scales illustrate that circuit optimization gains are not transferable from one configuration to another without empirical verification. Figure~\ref{fig:circuit_metrics} visualizes this comparison.

Second, alternative optimization frameworks such as ZX-calculus and Pytket proved useful for circuit-level analysis but were not practical drop-in optimizers for the parameterized model-training circuits used in this work. The end-to-end comparisons reported below therefore focus on the practically executable contrast between Qiskit L0 and Qiskit L3.

\begin{table*}[t]
\caption{Circuit metrics under L0 (baseline) and L3 (optimized) transpilation for each feature-map variant. Reduction percentages are relative to L0.}
\label{tab:circuit_metrics}
\centering
\setlength{\tabcolsep}{4.5pt}
\resizebox{0.98\textwidth}{!}{%
\begin{tabular}{|l|c|c|c|c|c|c|c|c|c|}
\hline
Feature map & Qubits & \multicolumn{2}{c|}{Depth} & \multicolumn{2}{c|}{Total gates} & \multicolumn{2}{c|}{CX gates} & Depth $\Delta$ (\%) & CX $\Delta$ (\%) \\
\cline{3-8}
 & & L0 & L3 & L0 & L3 & L0 & L3 & & \\
\hline
ZZ $r{=}2$ & 6 & 124 & 124 & 204 & 196 & 150 & 142 & 0.0 & $-$5.3 \\
ZZ $r{=}3$ & 6 & 179 & 181 & 306 & 323 & 225 & 200 & +1.1 & $-$11.1 \\
Pauli $r{=}2$ & 6 & 124 & 124 & 204 & 196 & 150 & 142 & 0.0 & $-$5.3 \\
ZZ $r{=}2$ & 10 & 237 & 195 & 604 & 592 & 474 & 372 & $-$17.7 & $-$21.5 \\
ZZ $r{=}3$ & 10 & 370 & 289 & 930 & 807 & 735 & 612 & $-$21.9 & $-$16.7 \\
Pauli $r{=}2$ & 10 & 237 & 195 & 604 & 592 & 474 & 372 & $-$17.7 & $-$21.5 \\
\hline
\end{tabular}%
}
\end{table*}

\Figure[t!](topskip=0pt,midskip=0pt,botskip=0pt)[width=0.94\textwidth]{figures/circuit_metrics_comparison.png}
{Circuit-level impact of L3 optimization. Reductions are modest at 6 qubits but substantial at 10 qubits, particularly for CX gate count.\label{fig:circuit_metrics}}

\subsection{Primary 6-Qubit Study}

The primary study provides the main inferential evidence because it is the only fully repeated factorial design completed on both datasets. Table~\ref{tab:primary6q_summary} summarizes the strongest observed configurations at the ideal and highest-noise endpoints. On IoT, the best ideal result was obtained by QWE with full entanglement and L3, reaching 95.62\% mean accuracy. Under the strongest noise level ($p_{1q}=0.05$), the strongest configuration was instead QVE with full entanglement, which reached 96.78\% mean accuracy. This non-monotonic behavior indicates that noise response depends on the model family rather than following a universal degradation trend. The largest IoT degradation was observed for QWE with full entanglement and L3, which dropped by 7.58 percentage points from ideal to the highest-noise condition, whereas QSVC with full entanglement dropped by approximately 5 percentage points. By contrast, QVE with full entanglement remained strong across the entire noise grid. Figure~\ref{fig:iot6q_noise} illustrates these 6-qubit IoT trajectories across model families, entanglement topologies, and optimization levels.

The UNSW dataset revealed an even sharper separation between robust and brittle configurations. Ideal accuracy was near ceiling for all models, with mean values around 99.6\%. At the highest noise level, full-entanglement QSVC collapsed to 9.53\% accuracy for both L0 and L3, and QWE with full entanglement collapsed to 9.53\% under L0. QWE with full entanglement and L3 partially recovered to 74.88\%, but remained far below its ideal performance. By contrast, QVE with full entanglement remained around 96.27\%, and all linear-entanglement models remained above 98.3\% at the same noise level. The central 6-qubit result on UNSW is therefore that model family and entanglement topology interact strongly with optimization under noise: full entanglement can be highly effective under ideal simulation yet catastrophically unstable for selected kernels under strong noise. Figure~\ref{fig:unsw6q_noise} makes this divergence clear by contrasting the collapse of full-entanglement QSVC and QWE with the stability of linear-entanglement and QVE configurations.

\Figure[t!](topskip=0pt,midskip=0pt,botskip=0pt)[width=0.94\textwidth]{figures/iot_6q_noise_response.png}
{IoT 6-qubit noise response. QSVC and QWE degrade more strongly under full entanglement, whereas QVE remains comparatively stable.\label{fig:iot6q_noise}}

\Figure[t!](topskip=0pt,midskip=0pt,botskip=0pt)[width=0.92\textwidth]{figures/unsw_6q_noise_response.png}
{UNSW 6-qubit noise response. Full entanglement is stronger in the ideal setting but more brittle under high noise, whereas linear entanglement remains robust.\label{fig:unsw6q_noise}}

\subsection{Optimization Effects Are Conditional, Not Uniform}

Holm-adjusted paired tests show that L3 is not universally superior to L0. Table~\ref{tab:holm_selected} lists representative statistically significant contrasts taken directly from the paired statistics files. On the IoT dataset, only 8 of the 60 paired 6-qubit comparisons were Holm-significant. The clearest L3 gains occurred for QWE with full entanglement at intermediate and high-noise settings, including a 2.57-point accuracy gain at $p_{1q}=0.02$. However, the same dataset also contains reversals: for QVE with full entanglement at $p_{1q}=0.02$, L0 exceeded L3 by 0.85 points. This is not a dataset where one can defensibly claim global L3 superiority.

The optimization effect is stronger on UNSW. There, 15 of the 60 paired comparisons were Holm-significant in the primary study. The largest effect appeared for QWE with full entanglement at $p_{1q}=0.05$, where L3 exceeded L0 by 65.35 accuracy points. Smaller but statistically significant L3 gains also appeared for noisy full-entanglement QSVC. Even so, these gains were regime-specific rather than universal. Many conditions showed negligible differences, and some significant contrasts reflected runtime rather than accuracy.

Figure~\ref{fig:heatmap_delta} presents a heatmap of the L3$-$L0 accuracy delta across all model-entanglement-noise combinations for both datasets. The IoT panel shows a scattered pattern with both positive and negative deltas, confirming that L3 is not uniformly beneficial. The UNSW panel reveals that most conditions show near-zero differences, with L3 gains concentrated exclusively in the QWE full-entanglement row at moderate-to-high noise levels, culminating in the extreme 65-point recovery at $p_{1q}=0.05$. These heatmaps complement the individual noise trajectories in Figs.~\ref{fig:iot6q_noise}--\ref{fig:unsw6q_noise} by showing the full factorial structure in a single view.

\Figure[t!](topskip=0pt,midskip=0pt,botskip=0pt)[width=0.94\textwidth]{figures/heatmap_accuracy_delta.png}
{Heatmap of L3$-$L0 accuracy delta (percentage points) across all 6-qubit conditions. Asterisks denote Holm-significant contrasts. IoT shows scattered small effects; UNSW shows concentrated gains for QWE/full at high noise.\label{fig:heatmap_delta}}

Figure~\ref{fig:cohens_d} provides a complementary view by plotting the paired Cohen's $d$ effect size for every condition. Most effects cluster near zero, falling below the conventional ``small'' threshold ($|d| < 0.2$). The few Holm-significant effects (marked as squares) appear at moderate-to-large $d$ values, predominantly for QWE/full on both datasets and QSVC/full on UNSW. The plot makes clear that statistically significant optimization gains are the exception rather than the rule.

\Figure[t!](topskip=0pt,midskip=0pt,botskip=0pt)[width=0.94\textwidth]{figures/cohens_d_effect_size.png}
{Paired Cohen's $d$ effect sizes for L3 vs.\ L0 across all 6-qubit conditions. Most effects are near zero; Holm-significant results (squares) appear only for selected model-noise combinations.\label{fig:cohens_d}}

\subsection{Supplementary 10-Qubit Extension}

The 10-qubit results are informative, but they should be interpreted as a supplementary extension rather than as a direct 6-qubit comparison because the model set, sample size, entanglement sweep, noise grid, and run count were all reduced to make the workload tractable. Table~\ref{tab:tenq_summary} collects the key robustness and runtime values for this extension. Within that restricted design, QVE again emerged as the most robust family. On IoT, QVE with full entanglement achieved 97.57\% ideal accuracy and still achieved 97.25\% at the highest noise level. QSVC, by contrast, fell from approximately 95.99\% to 90.67\%. On UNSW, QVE achieved 98.05\% in the ideal setting and remained between 96.52\% and 97.01\% at the highest noise level, whereas QSVC fell to 90.00\%.

The 10-qubit extension also highlights the computational cost of noisy higher-qubit kernels. Relative to ideal evaluation, total runtime increased by roughly 5.5$\times$ to 6.9$\times$ at the highest-noise setting. L0 versus L3 differences in this extension were mixed. IoT showed no Holm-significant accuracy advantage for L3. On UNSW, selected L3 gains did appear, including a 0.49-point improvement for QVE at $p_{1q}=0.05$, but QSVC exhibited unstable non-monotonic behavior across the reduced noise grid. Figure~\ref{fig:tenq_tradeoff} summarizes this extension as a joint accuracy-runtime trade-off rather than as a pure accuracy comparison. Taken together, these patterns reinforce the role of the 10-qubit block as a feasibility and robustness extension rather than a clean qubit-scaling experiment.

\Figure[t!](topskip=0pt,midskip=0pt,botskip=0pt)[width=0.94\textwidth]{figures/tenq_accuracy_runtime_tradeoff.png}
{Supplementary 10-qubit accuracy-runtime trade-off. QVE remains strong under noise, but 10-qubit evaluation incurs a sharp runtime increase.\label{fig:tenq_tradeoff}}

Figure~\ref{fig:boxplots} shows per-run accuracy distributions at four representative noise levels ($p_{1q} \in \{0, 0.005, 0.02, 0.05\}$) for both datasets. At the ideal and low-noise levels, all configurations produce tightly clustered distributions with similar medians. As noise increases, the IoT panel reveals increasing spread for QSVC and QWE with full entanglement, while QVE and linear-entanglement configurations maintain compact distributions. The UNSW panel at $p_{1q}=0.02$ shows the onset of collapse for full-entanglement QSVC, which falls to 9.53\%, while QWE with full entanglement remains above 98\% and all other configurations stay near ceiling. At $p_{1q}=0.05$, QWE full/L0 has also collapsed to 9.53\%, but QWE full/L3 shows partial recovery with substantial variance across runs (interquartile range spanning roughly 65--95\%), confirming that the 65-point L3 gain reported in Table~\ref{tab:holm_selected} is genuine but noisy.

\Figure[t!](topskip=0pt,midskip=0pt,botskip=0pt)[width=0.94\textwidth]{figures/boxplot_accuracy_distributions.png}
{Per-run accuracy distributions at four noise levels for both datasets. Boxes show L0 (blue) and L3 (red) distributions. Collapse and recovery patterns are visible at high noise for UNSW full-entanglement configurations.\label{fig:boxplots}}

\subsection{Cross-Validation}

Phase 5 cross-validation supports the ideal-simulation stability of the trained kernels. Table~\ref{tab:cv_results} reports the stratified 5-fold cross-validation metrics across all dataset and qubit configurations. For the primary 6-qubit study, mean cross-validated QSVC accuracy was approximately 95.4\% on IoT and 99.55\% on UNSW across the tested feature-map variants. In the 10-qubit extension, cross-validated QSVC accuracy remained around 96.3\% on IoT and 97.22\% on UNSW. ROC-AUC values exceed 0.96 in all configurations, confirming that the quantum kernels produce well-calibrated discriminative models. These results are consistent with the main noisy-study interpretation: ideal kernels can generalize well, but the ranking under noise depends strongly on circuit structure and model family.

\begin{table*}[t]
\caption{Stratified 5-fold cross-validation results for QSVC with different feature maps across all dataset and qubit configurations. Accuracy and F1 are reported as percentages.}
\label{tab:cv_results}
\centering
\setlength{\tabcolsep}{5pt}
\resizebox{0.98\textwidth}{!}{%
\begin{tabular}{|l|c|l|c|c|c|c|c|}
\hline
Dataset & Qubits & Feature map & Accuracy (\%) & F1 (\%) & MCC & ROC-AUC & Runs \\
\hline
IoT & 6 & ZZ & 95.41 $\pm$ 0.70 & 95.68 & 0.771 & 0.964 & 30 \\
IoT & 6 & Pauli & 95.49 $\pm$ 0.69 & 95.74 & 0.773 & 0.966 & 30 \\
IoT & 6 & ZZ Full & 95.43 $\pm$ 0.75 & 95.69 & 0.771 & 0.965 & 30 \\
IoT & 10 & ZZ & 96.26 $\pm$ 0.75 & 96.02 & 0.762 & 0.974 & 10 \\
IoT & 10 & Pauli & 96.27 $\pm$ 0.68 & 96.03 & 0.763 & 0.974 & 10 \\
IoT & 10 & ZZ Full & 96.28 $\pm$ 0.72 & 96.04 & 0.763 & 0.973 & 10 \\
UNSW & 6 & ZZ & 99.55 $\pm$ 0.20 & 99.55 & 0.974 & 0.997 & 30 \\
UNSW & 6 & Pauli & 99.55 $\pm$ 0.20 & 99.55 & 0.974 & 0.997 & 30 \\
UNSW & 6 & ZZ Full & 99.55 $\pm$ 0.20 & 99.55 & 0.974 & 0.997 & 30 \\
UNSW & 10 & ZZ & 97.22 $\pm$ 0.69 & 97.01 & 0.836 & 0.995 & 10 \\
UNSW & 10 & Pauli & 97.22 $\pm$ 0.69 & 97.01 & 0.836 & 0.995 & 10 \\
UNSW & 10 & ZZ Full & 97.22 $\pm$ 0.69 & 97.01 & 0.836 & 0.995 & 10 \\
\hline
\end{tabular}%
}
\end{table*}

\section{Discussion}
\label{sec:discussion}

The completed experiment leads to three main conclusions. First, circuit optimization matters, but not in the simplistic sense that lower depth automatically produces better classification. If that were the dominant mechanism, L3 would win uniformly and the most compressed circuits would always be the most robust. Instead, the results reveal a structured interaction in which optimization helps some models in some noise regimes, while entanglement choice and model architecture can dominate the outcome.

Second, the most robust design principle in this study is not maximal optimization but architectural diversity. QVE remained strong across both datasets and across both evidence tiers. This pattern suggests that heterogeneous quantum encodings provide a more reliable route to robustness than relying on a single expressive full-entanglement kernel, especially when that kernel is exposed to strong noise.

Third, the results call for caution when interpreting higher-qubit studies under realistic simulation budgets. The supplementary 10-qubit extension is scientifically useful because it shows that selected higher-qubit kernels, especially QVE, remain viable under noise. However, it does not support a direct claim that 10 qubits outperform 6 qubits, because the comparison is confounded by reduced sample size, reduced model set, reduced noise grid, restricted entanglement, and fewer runs. The appropriate interpretation is narrower: higher-qubit noisy kernel evaluation is feasible and informative, but substantially more expensive and therefore methodologically harder to study with the same breadth as the primary 6-qubit design.

From a deployment perspective, the main practical implication is selective rather than universal. Under noisy conditions, linear entanglement and QVE-style heterogeneous ensembles appear safer than highly entangled single-kernel models. L3 optimization should remain available as a design lever, but it should be justified empirically for the specific model and noise regime rather than assumed to be beneficial by default.

\section{Conclusion}
\label{sec:conclusion}

This work presented a five-phase study of circuit optimization for kernel-based quantum intrusion detection on two IoT cybersecurity datasets. The primary 6-qubit experiment showed that optimization effects are conditional: L3 improves robustness in selected noisy regimes, especially for full-entanglement QWE on UNSW, but it does not dominate L0 uniformly. Across both datasets, QVE was the most noise-robust model family, whereas full-entanglement QSVC and QWE could become brittle under strong noise. A reduced 10-qubit extension confirmed that selected higher-qubit kernels remain viable, but only at substantially greater computational cost and within a narrower experimental scope. Overall, the results support a practical design principle for near-term quantum kernel methods: robust performance depends jointly on optimization, entanglement topology, and model family, and cannot be inferred from circuit compression alone.

\section*{Appendix: Archived Result Files}

This appendix enumerates the complete set of 52 archived result files used to generate the tables, figures, and summary statements in the paper. The files are grouped by dataset and experiment batch so that readers can map each artifact to the evidence tier discussed in the main text.

\begin{table*}[t]
\caption{Archived result files for the IoT Original Distribution 6-qubit, 5000-sample batch (timestamp 20260310\_105202).}
\label{tab:archive_iot6}
\centering
\scriptsize
\setlength{\tabcolsep}{4pt}
\resizebox{0.98\textwidth}{!}{%
\begin{tabular}{|l|p{0.59\textwidth}|}
\hline
File & Contents \\
\hline
\texttt{circuit\_metrics\_iot\_original\_distribution\_6q\_5k\_v2\_noisy\_20260310\_105202.csv} & Circuit depth, total gates, CX gates, and feasibility diagnostics for the tested feature-map circuits. \\
\texttt{comprehensive\_model\_noise\_iot\_original\_distribution\_6q\_5k\_v2\_noisy\_20260310\_105202.csv} & Aggregated Phase 4 means across model family, entanglement, optimization level, and noise setting. \\
\texttt{comprehensive\_model\_noise\_runs\_iot\_original\_distribution\_6q\_5k\_v2\_noisy\_20260310\_105202.csv} & Per-run noisy-study outputs across the 30 paired seeds used in the primary study. \\
\texttt{comprehensive\_model\_noise\_stats\_iot\_original\_distribution\_6q\_5k\_v2\_noisy\_20260310\_105202.csv} & Paired L0 versus L3 statistics, Holm-adjusted tests, effect sizes, and confidence intervals for the noisy study. \\
\texttt{cv\_results\_iot\_original\_distribution\_6q\_5k\_v2\_noisy\_20260310\_105202.csv} & Stratified 5-fold cross-validation metrics for QSVC feature-map variants. \\
\texttt{experiment\_summary\_iot\_original\_distribution\_6q\_5k\_v2\_noisy\_20260310\_105202.json} & Batch-level configuration metadata, runtime totals, and artifact bookkeeping. \\
\texttt{journal\_summary\_iot\_original\_distribution\_6q\_5k\_v2\_noisy\_20260310\_105202.txt} & Human-readable narrative summary of the IoT 6-qubit batch for manuscript preparation. \\
\texttt{multi\_method\_model\_results\_iot\_original\_distribution\_6q\_5k\_v2\_noisy\_20260310\_105202.csv} & Model-level performance outputs for the supporting multi-method optimizer comparison. \\
\texttt{multi\_method\_optimization\_iot\_original\_distribution\_6q\_5k\_v2\_noisy\_20260310\_105202.csv} & Circuit-level comparison of Qiskit, PyZX, and Pytket optimization methods. \\
\texttt{multi\_method\_statistics\_iot\_original\_distribution\_6q\_5k\_v2\_noisy\_20260310\_105202.csv} & Statistical summaries supporting the multi-method optimization comparison. \\
\texttt{optimization\_comparison\_iot\_original\_distribution\_6q\_5k\_v2\_noisy\_20260310\_105202.csv} & Supporting unoptimized-versus-optimized model comparisons under ideal simulation. \\
\texttt{optimization\_levels\_iot\_original\_distribution\_6q\_5k\_v2\_noisy\_20260310\_105202.csv} & QSVC results across Qiskit transpiler optimization levels 0 through 3. \\
\texttt{optimization\_statistics\_iot\_original\_distribution\_6q\_5k\_v2\_noisy\_20260310\_105202.csv} & Statistical summaries for the optimization-level comparison phase. \\
\hline
\end{tabular}%
}
\end{table*}

\begin{table*}[t]
\caption{Archived result files for the UNSW 2018 IoT Botnet Final 10 Best 6-qubit, 5000-sample batch (timestamp 20260310\_105202).}
\label{tab:archive_unsw6}
\centering
\scriptsize
\setlength{\tabcolsep}{4pt}
\resizebox{0.98\textwidth}{!}{%
\begin{tabular}{|l|p{0.59\textwidth}|}
\hline
File & Contents \\
\hline
\texttt{circuit\_metrics\_unsw\_2018\_iot\_botnet\_final\_10\_best\_6q\_5k\_v2\_noisy\_20260310\_105202.csv} & Circuit depth, total gates, CX gates, and feasibility diagnostics for the tested feature-map circuits. \\
\texttt{comprehensive\_model\_noise\_unsw\_2018\_iot\_botnet\_final\_10\_best\_6q\_5k\_v2\_noisy\_20260310\_105202.csv} & Aggregated Phase 4 means across model family, entanglement, optimization level, and noise setting. \\
\texttt{comprehensive\_model\_noise\_runs\_unsw\_2018\_iot\_botnet\_final\_10\_best\_6q\_5k\_v2\_noisy\_20260310\_105202.csv} & Per-run noisy-study outputs across the 30 paired seeds used in the primary study. \\
\texttt{comprehensive\_model\_noise\_stats\_unsw\_2018\_iot\_botnet\_final\_10\_best\_6q\_5k\_v2\_noisy\_20260310\_105202.csv} & Paired L0 versus L3 statistics, Holm-adjusted tests, effect sizes, and confidence intervals for the noisy study. \\
\texttt{cv\_results\_unsw\_2018\_iot\_botnet\_final\_10\_best\_6q\_5k\_v2\_noisy\_20260310\_105202.csv} & Stratified 5-fold cross-validation metrics for QSVC feature-map variants. \\
\texttt{experiment\_summary\_unsw\_2018\_iot\_botnet\_final\_10\_best\_6q\_5k\_v2\_noisy\_20260310\_105202.json} & Batch-level configuration metadata, runtime totals, and artifact bookkeeping. \\
\texttt{journal\_summary\_unsw\_2018\_iot\_botnet\_final\_10\_best\_6q\_5k\_v2\_noisy\_20260310\_105202.txt} & Human-readable narrative summary of the UNSW 6-qubit batch for manuscript preparation. \\
\texttt{multi\_method\_model\_results\_unsw\_2018\_iot\_botnet\_final\_10\_best\_6q\_5k\_v2\_noisy\_20260310\_105202.csv} & Model-level performance outputs for the supporting multi-method optimizer comparison. \\
\texttt{multi\_method\_optimization\_unsw\_2018\_iot\_botnet\_final\_10\_best\_6q\_5k\_v2\_noisy\_20260310\_105202.csv} & Circuit-level comparison of Qiskit, PyZX, and Pytket optimization methods. \\
\texttt{multi\_method\_statistics\_unsw\_2018\_iot\_botnet\_final\_10\_best\_6q\_5k\_v2\_noisy\_20260310\_105202.csv} & Statistical summaries supporting the multi-method optimization comparison. \\
\texttt{optimization\_comparison\_unsw\_2018\_iot\_botnet\_final\_10\_best\_6q\_5k\_v2\_noisy\_20260310\_105202.csv} & Supporting unoptimized-versus-optimized model comparisons under ideal simulation. \\
\texttt{optimization\_levels\_unsw\_2018\_iot\_botnet\_final\_10\_best\_6q\_5k\_v2\_noisy\_20260310\_105202.csv} & QSVC results across Qiskit transpiler optimization levels 0 through 3. \\
\texttt{optimization\_statistics\_unsw\_2018\_iot\_botnet\_final\_10\_best\_6q\_5k\_v2\_noisy\_20260310\_105202.csv} & Statistical summaries for the optimization-level comparison phase. \\
\hline
\end{tabular}%
}
\end{table*}

\begin{table*}[t]
\caption{Archived result files for the IoT Original Distribution 10-qubit, 2500-sample batch (timestamp 20260323\_114741).}
\label{tab:archive_iot10}
\centering
\scriptsize
\setlength{\tabcolsep}{4pt}
\resizebox{0.98\textwidth}{!}{%
\begin{tabular}{|l|p{0.59\textwidth}|}
\hline
File & Contents \\
\hline
\texttt{circuit\_metrics\_iot\_original\_distribution\_10q\_2p5k\_v2\_noisy\_20260323\_114741.csv} & Circuit depth, total gates, CX gates, and feasibility diagnostics for the tested feature-map circuits. \\
\texttt{comprehensive\_model\_noise\_iot\_original\_distribution\_10q\_2p5k\_v2\_noisy\_20260323\_114741.csv} & Aggregated Phase 4 means across model family, optimization level, and reduced noise grid for the 10-qubit extension. \\
\texttt{comprehensive\_model\_noise\_runs\_iot\_original\_distribution\_10q\_2p5k\_v2\_noisy\_20260323\_114741.csv} & Per-run noisy-study outputs across the 10 paired seeds used in the 10-qubit extension. \\
\texttt{comprehensive\_model\_noise\_stats\_iot\_original\_distribution\_10q\_2p5k\_v2\_noisy\_20260323\_114741.csv} & Paired L0 versus L3 statistics, Holm-adjusted tests, effect sizes, and confidence intervals for the 10-qubit noisy study. \\
\texttt{cv\_results\_iot\_original\_distribution\_10q\_2p5k\_v2\_noisy\_20260323\_114741.csv} & Stratified 5-fold cross-validation metrics for QSVC feature-map variants at 10 qubits. \\
\texttt{experiment\_summary\_iot\_original\_distribution\_10q\_2p5k\_v2\_noisy\_20260323\_114741.json} & Batch-level configuration metadata, runtime totals, and artifact bookkeeping. \\
\texttt{journal\_summary\_iot\_original\_distribution\_10q\_2p5k\_v2\_noisy\_20260323\_114741.txt} & Human-readable narrative summary of the IoT 10-qubit batch for manuscript preparation. \\
\texttt{multi\_method\_model\_results\_iot\_original\_distribution\_10q\_2p5k\_v2\_noisy\_20260323\_114741.csv} & Model-level performance outputs for the supporting multi-method optimizer comparison. \\
\texttt{multi\_method\_optimization\_iot\_original\_distribution\_10q\_2p5k\_v2\_noisy\_20260323\_114741.csv} & Circuit-level comparison of Qiskit, PyZX, and Pytket optimization methods. \\
\texttt{multi\_method\_statistics\_iot\_original\_distribution\_10q\_2p5k\_v2\_noisy\_20260323\_114741.csv} & Statistical summaries supporting the multi-method optimization comparison. \\
\texttt{optimization\_comparison\_iot\_original\_distribution\_10q\_2p5k\_v2\_noisy\_20260323\_114741.csv} & Supporting unoptimized-versus-optimized model comparisons under ideal simulation. \\
\texttt{optimization\_levels\_iot\_original\_distribution\_10q\_2p5k\_v2\_noisy\_20260323\_114741.csv} & QSVC results across Qiskit transpiler optimization levels 0 through 3. \\
\texttt{optimization\_statistics\_iot\_original\_distribution\_10q\_2p5k\_v2\_noisy\_20260323\_114741.csv} & Statistical summaries for the optimization-level comparison phase. \\
\hline
\end{tabular}%
}
\end{table*}

\begin{table*}[t]
\caption{Archived result files for the UNSW 2018 IoT Botnet Final 10 Best 10-qubit, 2500-sample batch (timestamp 20260323\_114741).}
\label{tab:archive_unsw10}
\centering
\scriptsize
\setlength{\tabcolsep}{4pt}
\resizebox{0.98\textwidth}{!}{%
\begin{tabular}{|l|p{0.59\textwidth}|}
\hline
File & Contents \\
\hline
\texttt{circuit\_metrics\_unsw\_2018\_iot\_botnet\_final\_10\_best\_10q\_2p5k\_v2\_noisy\_20260323\_114741.csv} & Circuit depth, total gates, CX gates, and feasibility diagnostics for the tested feature-map circuits. \\
\texttt{comprehensive\_model\_noise\_unsw\_2018\_iot\_botnet\_final\_10\_best\_10q\_2p5k\_v2\_noisy\_20260323\_114741.csv} & Aggregated Phase 4 means across model family, optimization level, and reduced noise grid for the 10-qubit extension. \\
\texttt{comprehensive\_model\_noise\_runs\_unsw\_2018\_iot\_botnet\_final\_10\_best\_10q\_2p5k\_v2\_noisy\_20260323\_114741.csv} & Per-run noisy-study outputs across the 10 paired seeds used in the 10-qubit extension. \\
\texttt{comprehensive\_model\_noise\_stats\_unsw\_2018\_iot\_botnet\_final\_10\_best\_10q\_2p5k\_v2\_noisy\_20260323\_114741.csv} & Paired L0 versus L3 statistics, Holm-adjusted tests, effect sizes, and confidence intervals for the 10-qubit noisy study. \\
\texttt{cv\_results\_unsw\_2018\_iot\_botnet\_final\_10\_best\_10q\_2p5k\_v2\_noisy\_20260323\_114741.csv} & Stratified 5-fold cross-validation metrics for QSVC feature-map variants at 10 qubits. \\
\texttt{experiment\_summary\_unsw\_2018\_iot\_botnet\_final\_10\_best\_10q\_2p5k\_v2\_noisy\_20260323\_114741.json} & Batch-level configuration metadata, runtime totals, and artifact bookkeeping. \\
\texttt{journal\_summary\_unsw\_2018\_iot\_botnet\_final\_10\_best\_10q\_2p5k\_v2\_noisy\_20260323\_114741.txt} & Human-readable narrative summary of the UNSW 10-qubit batch for manuscript preparation. \\
\texttt{multi\_method\_model\_results\_unsw\_2018\_iot\_botnet\_final\_10\_best\_10q\_2p5k\_v2\_noisy\_20260323\_114741.csv} & Model-level performance outputs for the supporting multi-method optimizer comparison. \\
\texttt{multi\_method\_optimization\_unsw\_2018\_iot\_botnet\_final\_10\_best\_10q\_2p5k\_v2\_noisy\_20260323\_114741.csv} & Circuit-level comparison of Qiskit, PyZX, and Pytket optimization methods. \\
\texttt{multi\_method\_statistics\_unsw\_2018\_iot\_botnet\_final\_10\_best\_10q\_2p5k\_v2\_noisy\_20260323\_114741.csv} & Statistical summaries supporting the multi-method optimization comparison. \\
\texttt{optimization\_comparison\_unsw\_2018\_iot\_botnet\_final\_10\_best\_10q\_2p5k\_v2\_noisy\_20260323\_114741.csv} & Supporting unoptimized-versus-optimized model comparisons under ideal simulation. \\
\texttt{optimization\_levels\_unsw\_2018\_iot\_botnet\_final\_10\_best\_10q\_2p5k\_v2\_noisy\_20260323\_114741.csv} & QSVC results across Qiskit transpiler optimization levels 0 through 3. \\
\texttt{optimization\_statistics\_unsw\_2018\_iot\_botnet\_final\_10\_best\_10q\_2p5k\_v2\_noisy\_20260323\_114741.csv} & Statistical summaries for the optimization-level comparison phase. \\
\hline
\end{tabular}%
}
\end{table*}

\section*{Acknowledgment}
% TODO: Add acknowledgment text. Use singular heading (no trailing 's').
% Spell as ``acknowledgment'' (American English, no 'e' after 'g').


\begin{thebibliography}{00}

\bibitem{ullah2020scheme} I. Ullah and Q. H. Mahmoud, ``A scheme for generating a dataset for anomalous activity detection in IoT networks,'' in \emph{Proc. 33rd Can. Conf. Artif. Intell. (Canadian AI)}, Ottawa, ON, Canada, 2020, pp. 508--520, DOI: 10.1007/978-3-030-47358-7\_52.

%% NISQ Era and Quantum Hardware
\bibitem{Preskill2018} J. Preskill, ``Quantum computing in the NISQ era and beyond,'' \emph{Quantum}, vol. 2, p. 79, 2018, DOI: 10.22331/q-2018-08-06-79.

\bibitem{Cross2019} A. W. Cross, L. S. Bishop, S. Sheldon, P. D. Nation, and J. M. Gambetta, ``Validating quantum computers using randomized model circuits,'' \emph{Phys. Rev. A}, vol. 100, no. 3, p. 032328, 2019, DOI: 10.1103/PhysRevA.100.032328.

%% Compilers and Transpilers
\bibitem{Qiskit2024} {Qiskit Contributors}, ``Qiskit: An open-source framework for quantum computing,'' 2024. [Online]. Available: https://qiskit.org

\bibitem{Cirq2021} {Cirq Developers}, ``Cirq: A Python framework for creating, editing, and invoking noisy intermediate-scale quantum circuits,'' 2021. [Online]. Available: https://quantumai.google/cirq

\bibitem{Sivarajah2020} S. Sivarajah, S. Dilkes, A. Cowtan, W. Simmons, A. Edgington, and R. Duncan, ``t$|$ket$\rangle$: A retargetable compiler for NISQ devices,'' \emph{Quantum Sci. Technol.}, vol. 6, no. 1, p. 014003, 2021, DOI: 10.1088/2058-9565/ab8e92.

\bibitem{Cross2022} A. W. Cross \emph{et al.}, ``OpenQASM 3: A broader and deeper quantum assembly language,'' \emph{ACM Trans. Quantum Comput.}, vol. 3, no. 3, pp. 1--50, 2022, DOI: 10.1145/3505636.

\bibitem{Dilillo2023} N. Dilillo, E. Giusto, E. Dri, B. Baheri \emph{et al.}, ``Understanding the effect of transpilation in the reliability of quantum circuits,'' in \emph{Proc. IEEE Int. Symp. Defect Fault Tolerance VLSI Nanotechnol. Syst. (DFT)}, 2023, DOI: 10.1109/DFT59622.2023.10313838.

\bibitem{Nation2025} P. D. Nation, A. A. Saki, S. Brandhofer, L. Bello \emph{et al.}, ``Benchmarking the performance of quantum computing software for quantum circuit creation, manipulation and compilation,'' \emph{Nature Comput. Sci.}, 2025, DOI: 10.1038/s43588-025-00792-y.

\bibitem{Salm2021} M. Salm, J. Barzen, F. Leymann, and B. Weder, ``Automating the comparison of quantum compilers for quantum circuits,'' in \emph{Symp. Summer School Service-Oriented Comput.}, Springer, 2021, pp. 64--80, DOI: 10.1007/978-3-030-87568-8\_4.

\bibitem{Nam2018} Y. Nam, N. J. Ross, Y. Su, A. M. Childs, and D. Maslov, ``Automated optimization of large quantum circuits with continuous parameters,'' \emph{npj Quantum Inf.}, vol. 4, no. 1, p. 23, 2018, DOI: 10.1038/s41534-018-0072-4.

\bibitem{Liu2021} J. Liu, L. Bello, and H. Zhou, ``Relaxed peephole optimization:\ A novel compiler optimization for quantum circuits,'' in \emph{Proc. IEEE/ACM Int. Symp. Code Generation Optim. (CGO)}, 2021, DOI: 10.1109/CGO51591.2021.9370310.

%% ZX-Calculus
\bibitem{Coecke2011} B. Coecke and R. Duncan, ``Interacting quantum observables: Categorical algebra and diagrammatics,'' \emph{New J. Phys.}, vol. 13, no. 4, p. 043016, 2011, DOI: 10.1088/1367-2630/13/4/043016.

\bibitem{vandeWetering2020} J. van de Wetering, ``ZX-calculus for the working quantum computer scientist,'' \emph{arXiv:2012.13966}, 2020.

\bibitem{Kissinger2020} A. Kissinger and J. van de Wetering, ``Reducing the number of non-Clifford gates in quantum circuits,'' \emph{Phys. Rev. A}, vol. 102, no. 2, p. 022406, 2020, DOI: 10.1103/PhysRevA.102.022406.

\bibitem{Kissinger2019} A. Kissinger and J. van de Wetering, ``PyZX: Large scale automated diagrammatic reasoning,'' in \emph{Electron. Proc. Theor. Comput. Sci.}, vol. 318, pp. 229--241, 2020, DOI: 10.4204/EPTCS.318.14.

\bibitem{deBeaudrap2020} N. de Beaudrap, X. Bian, and Q. Wang, ``Fast and effective techniques for T-count reduction via spider nest identities,'' \emph{arXiv:2004.05164}, 2020.

\bibitem{Fischbach2025} T. Fischbach, P. Talbot, and P. Bouvry, ``A review on quantum circuit optimization using ZX-calculus,'' \emph{arXiv:2509.20663}, 2025.

\bibitem{vandeWetering2025} J. van de Wetering, R. Yeung, T. Laakkonen \emph{et al.}, ``Optimal compilation of parametrised quantum circuits,'' \emph{Quantum}, vol. 9, p. 1828, 2025, DOI: 10.22331/q-2025-08-27-1828.

%% Synthesis and Compilation Survey
\bibitem{Yan2024} G. Yan, W. Wu, Y. Chen, K. Pan, X. Lu, Z. Zhou \emph{et al.}, ``Quantum circuit synthesis and compilation optimization: Overview and prospects,'' \emph{arXiv:2407.00736}, 2024.

%% ML for Circuit Optimization
\bibitem{Fosel2021} T. F{\"o}sel, M. Y. Niu, F. Marquardt, and L. Li, ``Quantum circuit optimization with deep reinforcement learning,'' \emph{arXiv:2103.07585}, 2021.

\bibitem{Paler2023} A. Paler, L. M. Sasu, A.-C. Florea, and R. Andonie, ``Machine learning optimization of quantum circuit layouts,'' \emph{ACM Trans. Quantum Comput.}, vol. 4, no. 1, pp. 1--25, 2023, DOI: 10.1145/3565271.

\bibitem{Ruiz2025} F. J. R. Ruiz \emph{et al.}, ``Quantum circuit optimization with AlphaTensor,'' \emph{Nature Mach. Intell.}, 2025, DOI: 10.1038/s42256-025-01001-1.

\bibitem{Kremer2024} D. Kremer \emph{et al.}, ``Practical and efficient quantum circuit synthesis and transpiling with reinforcement learning,'' \emph{arXiv:2405.13196}, 2024.

\bibitem{Wang2022QuantumNAS} H. Wang \emph{et al.}, ``QuantumNAS: Noise-adaptive search for robust quantum circuits,'' in \emph{Proc. IEEE Int. Symp. High-Perform. Comput. Architecture (HPCA)}, 2022, DOI: 10.1109/HPCA53966.2022.00065.

%% Quantum Kernel Methods and QML
\bibitem{Schuld2019} M. Schuld and N. Killoran, ``Quantum machine learning in feature Hilbert spaces,'' \emph{Phys. Rev. Lett.}, vol. 122, no. 4, p. 040504, 2019, DOI: 10.1103/PhysRevLett.122.040504.

\bibitem{Havlicek2019} V. Havl\'{i}\v{c}ek \emph{et al.}, ``Supervised learning with quantum-enhanced feature spaces,'' \emph{Nature}, vol. 567, no. 7747, pp. 209--212, 2019, DOI: 10.1038/s41586-019-0980-2.

\bibitem{Huang2021} H.-Y. Huang \emph{et al.}, ``Power of data in quantum machine learning,'' \emph{Nature Commun.}, vol. 12, no. 1, p. 2631, 2021, DOI: 10.1038/s41467-021-22539-9.

\bibitem{Sim2019} S. Sim, P. D. Johnson, and A. Aspuru-Guzik, ``Expressibility and entangling capability of parameterized quantum circuits for hybrid quantum-classical algorithms,'' \emph{Adv. Quantum Technol.}, vol. 2, no. 12, p. 1900070, 2019, DOI: 10.1002/qute.201900070.

\bibitem{McClean2018} J. R. McClean, S. Boixo, V. N. Smelyanskiy, R. Babbush, and H. Neven, ``Barren plateaus in quantum neural network training landscapes,'' \emph{Nature Commun.}, vol. 9, no. 1, p. 4812, 2018, DOI: 10.1038/s41467-018-07090-4.

\bibitem{Holmes2022} Z. Holmes, K. Sharma, M. Cerezo, and P. J. Coles, ``Connecting ansatz expressibility to gradient magnitudes and barren plateaus,'' \emph{PRX Quantum}, vol. 3, no. 1, p. 010313, 2022, DOI: 10.1103/PRXQuantum.3.010313.

%% Quantum Ensemble Methods
\bibitem{Schuld2018Ensemble} M. Schuld and F. Petruccione, ``Quantum ensembles of quantum classifiers,'' \emph{Sci. Rep.}, vol. 8, p. 2772, 2018, DOI: 10.1038/s41598-018-20403-3.

\bibitem{Qin2022} R. Qin, Z. Liang, J. Cheng, P. Kogge, and Y. Shi, ``Improving quantum classifier performance in NISQ computers by voting strategy from ensemble learning,'' \emph{arXiv:2210.01656}, 2022.

\bibitem{Tolotti2023} E. Tolotti, E. Zardini, E. Blanzieri, and D. Pastorello, ``Ensembles of quantum classifiers,'' \emph{arXiv:2311.09750}, 2023.

\bibitem{Incudini2023} M. Incudini, M. Grossi, A. Ceschini, A. Mandarino \emph{et al.}, ``Resource saving via ensemble techniques for quantum neural networks,'' \emph{Quantum Mach. Intell.}, vol. 5, p. 2, 2023, DOI: 10.1007/s42484-023-00126-z.

\bibitem{Li2024ensemble} Q. Li, Y. Huang, X. Hou, Y. Li, X. Wang, and A. Bayat, ``Ensemble-learning error mitigation for variational quantum shallow-circuit classifiers,'' \emph{Phys. Rev. Research}, vol. 6, p. 013027, 2024, DOI: 10.1103/PhysRevResearch.6.013027.

%% Intrusion Detection and Cybersecurity
\bibitem{Buczak2016} A. L. Buczak and E. Guven, ``A survey of data mining and machine learning methods for cyber security intrusion detection,'' \emph{IEEE Commun. Surveys Tuts.}, vol. 18, no. 2, pp. 1153--1176, 2016, DOI: 10.1109/COMST.2015.2494502.

\bibitem{Aldweesh2020} A. Aldweesh, A. Derhab, and A. Z. Emam, ``Deep learning approaches for anomaly-based intrusion detection systems: A survey, taxonomy, and open issues,'' \emph{Knowl.-Based Syst.}, vol. 189, p. 105124, 2020, DOI: 10.1016/j.knosys.2019.105124.

\bibitem{Vinayakumar2019} R. Vinayakumar \emph{et al.}, ``Deep learning approach for intelligent intrusion detection system,'' \emph{IEEE Access}, vol. 7, pp. 41525--41550, 2019, DOI: 10.1109/ACCESS.2019.2895334.

\bibitem{Cook2019} A. A. Cook, G. M{\i}s{\i}rl{\i}, and Z. Fan, ``Anomaly detection for IoT time-series data: A survey,'' \emph{IEEE Internet Things J.}, vol. 7, no. 7, pp. 6481--6494, 2019, DOI: 10.1109/JIOT.2019.2958185.

%% Quantum for Cybersecurity and IoT
\bibitem{Kalinin2023} M. Kalinin and V. Krundyshev, ``Security intrusion detection using quantum machine learning techniques,'' \emph{J. Comput. Virol. Hacking Techn.}, vol. 19, pp. 125--136, 2023, DOI: 10.1007/s11416-022-00435-0.

\bibitem{Kumar2025} R. Kumar and M. Swarnkar, ``QuIDS: A quantum support vector machine-based intrusion detection system for IoT networks,'' \emph{J. Netw. Comput. Appl.}, vol. 234, p. 104092, 2025, DOI: 10.1016/j.jnca.2024.104092.

\bibitem{Alomari2023} A. Alomari and S. A. P. Kumar, ``DEQSVC: Dimensionality reduction and encoding technique for quantum support vector classifier approach to detect DDoS attacks,'' \emph{IEEE Access}, vol. 11, pp. 113207--113220, 2023, DOI: 10.1109/ACCESS.2023.3273991.

\bibitem{Elsedimy2024} E. I. Elsedimy, H. Elhadidy, and S. M. M. Abohashish, ``A novel intrusion detection system based on a hybrid quantum support vector machine and improved Grey Wolf optimizer,'' \emph{Cluster Comput.}, vol. 27, pp. 7269--7291, 2024, DOI: 10.1007/s10586-024-04458-8.

\bibitem{Cultice2025} T. Cultice, M. S. H. Onim, A. Giani, and H. Thapliyal, ``Quantum-hybrid support vector machines for anomaly detection in industrial control systems,'' \emph{arXiv:2506.17824}, 2025.

\bibitem{Houkan2024} A. Houkan, A. K. Sahoo, S. P. Gochhayat, and P. K. Sahoo, ``Enhancing security in industrial IoT networks: Machine learning solutions for feature selection and reduction,'' \emph{IEEE Access}, vol. 12, pp. 170476--170493, 2024, DOI: 10.1109/ACCESS.2024.3489898.

\bibitem{Adefemi2025} K. O. Adefemi, M. B. Mutanga, and O. A. Alimi, ``A hybrid CNN--GRU deep learning model for IoT network intrusion detection,'' \emph{J. Sensor Actuator Netw.}, vol. 14, no. 5, p. 96, 2025.

%% Hardware Optimization and Noise Studies
\bibitem{Murali2020} P. Murali, J. M. Baker, A. Javadi-Abhari, F. T. Chong, and M. Martonosi, ``Software mitigation of crosstalk on noisy intermediate-scale quantum computers,'' in \emph{Proc. 25th Int. Conf. Archit. Support Program. Lang. Oper. Syst. (ASPLOS)}, 2020, pp. 1001--1016, DOI: 10.1145/3373376.3378477.

%% Authors' Prior Work
% TODO: Replace author names with actual names once published
\bibitem{AuthorQVEQWE2025} [Author names], ``Novel quantum ensemble machine learning models for IoT intrusion detection,'' submitted to \emph{IEEE Dallas Circuits Syst. Conf. (DCAS)}, 2025.

\end{thebibliography}

% TODO: Add author biographies. Each biography should include:
%   - First paragraph: place/date of birth, education (degrees, fields, institutions, years)
%   - Second paragraph: work experience (use pronouns he/she, not last name)
%   - Third paragraph: title + last name, professional society memberships, awards
% Author photos should be 1 inch wide by 1.25 inches tall.
%\begin{IEEEbiography}[{\includegraphics[width=1in,height=1.25in,clip,keepaspectratio]{author1.png}}]{Author Name}
% Biography text here.
%\end{IEEEbiography}

\EOD

\end{document}
